\documentclass[10pt,letterpaper]{article}
\usepackage{cloudsuite}
\renewcommand{\DocShort}{CH. 01}
\begin{document}
\coverpage{Chapter 1: Introduction}{Cornell study notes and applied review | ORIENTATION}{Concepts \textbullet{} decisions \textbullet{} verification artifacts}
\sourceStatement
\tableofcontents
\newpage
\section{Learning contract}
\begin{takeaway}
Define the boundary between vendor-neutral concepts and vendor-specific implementation.
\end{takeaway}
\textbf{Chapter scope:} Scope, audience, organization and references.\\
\sourceRef{1}{1}
\subsection{Mapped source sections}
\begin{itemize}
\item 1.1 Objectives of This Book
\item 1.2 What This Book Does Not Cover
\item 1.3 Who This Book Is For
\item 1.4 How This Book Is Organized
\item 1.5 Resources
\end{itemize}
\section{Cornell cue and notes matrix}
\small\setlength{\LTleft}{0pt}\setlength{\LTright}{0pt}
\begin{longtable}{>{\RaggedRight\arraybackslash}p{0.29\textwidth}>{\RaggedRight\arraybackslash}p{0.63\textwidth}}
\toprule\textbf{Cue / essential question}&\textbf{Explanatory notes and reasoning}\\\midrule\endfirsthead\toprule\textbf{Cue}&\textbf{Notes}\\\midrule\endhead
\textbf{Why vendor-neutral terminology?} & A conceptual vocabulary separates enduring platform capabilities from product names. Record resource, responsibility and trust assumptions before choosing implementation-specific services.\\[3pt]\textcolor{petrol}{\emph{Practice:} Write a one-page problem statement that names a cloud consumer, an expected service and an observable outcome.}\\[8pt]\midrule
\textbf{What does the book deliberately exclude?} & The source emphasizes broadly applicable concepts rather than detailed product implementation. Product-specific API commands and prices therefore belong in a separate implementation appendix, not as claims about the source.\\[3pt]\textcolor{petrol}{\emph{Practice:} Flag each proposed command or cloud product as external material.}\\[8pt]\midrule
\textbf{How is the book structured?} & The learning sequence proceeds from foundational concepts to mechanisms, assembled architectures and operational considerations. Earlier vocabulary makes later architecture comparisons more precise.\\[3pt]\textcolor{petrol}{\emph{Practice:} Draw the dependency chain of the five parts.}\\[8pt]\midrule
\textbf{Who benefits from the taxonomy?} & Architects, engineering practitioners, decision-makers and students need a shared language for technical and business tradeoffs.\\[3pt]\textcolor{petrol}{\emph{Practice:} Explain a cloud design in a sentence each to an engineer and a budget owner.}\\[8pt]\midrule
\textbf{What counts as a substantiated design decision?} & A decision ties requirements to a named concept, an observable constraint and a test. An architecture label alone does not demonstrate reliability, security or value.\\[3pt]\textcolor{petrol}{\emph{Practice:} Start a design decision record with assumptions and acceptance evidence.}\\[8pt]\midrule
\end{longtable}\normalsize
\section{Synthesis and architecture connections}
Translate the chapter vocabulary into one explicit service requirement, one boundary, one owner and one measurable test. Later mechanism and architecture chapters provide possible realizations; the requirements remain the acceptance baseline.
\begin{workbox}[title={Applied deliverable}]
Define the boundary between vendor-neutral concepts and vendor-specific implementation.\\[4pt]Create a one-page architecture or operations artifact with: (1) assumptions, (2) a decision and alternatives, (3) control ownership, (4) observable acceptance criteria, and (5) a failed-case test.
\end{workbox}
\section{Active recall}
\begin{enumerate}
\item Why vendor-neutral terminology?
\item What does the book deliberately exclude?
\item How is the book structured?
\item Who benefits from the taxonomy?
\end{enumerate}
\subsection{Answer-check rubric}
\begin{itemize}
\item A correct response distinguishes the named concepts rather than relying on vendor names.
\item The explanation identifies an operational consequence, boundary or tradeoff.
\item At least one measurement, test or documentary artifact could demonstrate the claim.
\end{itemize}
\section{Blank Cornell capture sheet}
\textbf{Topic / source subsection:}\enspace\rule{0.65\linewidth}{0.4pt}\\[6pt]
\begin{longtable}{p{0.28\textwidth}p{0.64\textwidth}}
\toprule Cues / questions & Notes / evidence \\ \midrule
\rule{0pt}{42pt} & \\ \midrule
\rule{0pt}{42pt} & \\ \midrule
\rule{0pt}{42pt} & \\ \midrule
\rule{0pt}{42pt} & \\ \midrule
\rule{0pt}{42pt} & \\ \midrule
\end{longtable}
\textbf{Summary in three sentences:}\\[16pt]\hrulefill\\[16pt]\hrulefill\\[16pt]\hrulefill

\refs
\end{document}
